\documentclass[11pt]{article}

\usepackage[margin=1.5in]{geometry}
\usepackage{lecnotes}
% \usepackage{mpass}
% \lstset{style=verb,language=mpass}
\input{lmacros}

\newcommand{\course}{15-316: Software Foundations of Security \& Privacy}
\newcommand{\lecturer}{Matt Fredrikson}
\newcommand{\lecdate}{March 10, 2026}
\newcommand{\lecnum}{14}
\newcommand{\lectitle}{Termination-Sensitive Noninterference}
\newcommand{\courseurl}{https://15316-cmu.github.io/2026/}

\begin{document}

\maketitle

\section{Introduction}

Our assumption so far in the semantic definition of information flow has been
that an attacker can only compare the values of low-security variables in the
poststates.
\begin{quote}
  We define $\Sigma \models \alpha\; \ms{secure}$ (program $\alpha$ satisfies
  \emph{termination-insensitive noninterference with respect to policy
    $\Sigma$}) \newline iff
  \begin{enumerate}[]
  \item for all $\ell$, $\omega_1$, $\omega_2$, $\nu_1$, and $\nu_2$,
  \item whenever $\Sigma \vdash \omega_1 \approx_\ell \omega_2$,
  \item and $\omega_1 \lbb \alpha \rbb \nu_1$,
  \item and $\omega_2 \lbb \alpha \rbb \nu_2$
  \item then $\Sigma \vdash \nu_1 \approx_\ell \nu_2$.
  \end{enumerate}
\end{quote}
This ignores that possibility that an attacker might notice that a program
does not terminate.  For example, with $(x : \ms{H})$ and $\ms{L} \sqsubset \ms{H}$
\begin{tabbing}
  \qquad $\mb{if}\; x > 5\; \mb{then}\; (\mb{while}\; \top\; \mb{skip})\; \mb{else}\; \mb{skip}$
\end{tabbing}
an observer can tell if $x > 5$ based on whether the program terminates.  But
according to our policy it is secure at level $\ms{L}$!  In this example, that's
easy to see because whenever $\nu_1$ and $\nu_2$ both exist, then
$\nu_1 \approx_{\ms{L}} \nu_2$ because there aren't even any low-security
variables.  If $\alpha$ does not terminate (that is, has no poststate) in one or
both of the two initial states, then the implication is vacuously true.

The first goal of today's lecture is to sharpen our definition of
noninterference so that nontermination is considered observable, and the example
above would be rejected.  The second goal will be to revise the information flow
type system so it is sound with respect to the sharpened definition.  For
general expositions of today's topic, see~\citet{Hedin12sss,Sabelfeld03comm},
and earlier seminal work by \citet{Volpano97csfw}.

This lecture is also a preparation for the next lecture where we talk about
\emph{timing attacks} where information can leak by observing how long a program
takes to execute.  These kind of attacks are quite common and real concerns in
actual systems \citep{Brumley05cn}, so it is important to study them from the
programming perspective.

\section{Termination-Sensitive Noninterference}

Intuitively, we'd like to say that if $\omega_1 \lbb \alpha \rbb \nu_1$ and
$\omega_2 \lbb \alpha \rbb \nu_2$ both hold, then the poststates still must be
equivalent to an observer at level $\ell$.  Moreover, if $\alpha$ has no
poststate starting from $\omega_1$, then it should also have no poststate
starting from $\omega_2$.  In this latter case, there are no poststates
to compare.  The traditional definition \citep{Volpano97csfw} captures this by stating that
if $\omega_1 \lbb \alpha \rbb \nu_1$ then there must also be a $\nu_2$ such
that $\omega_2 \lbb \alpha \rbb \nu_2$, and the two poststates must be
observably equivalent.  This slightly asymmetric definition is correct due to
the symmetry of the $\approx_\ell$ relation.  If $\alpha$ terminates in prestate
$\omega_1$ but does not in $\omega_2$, then we can just swap the two prestates
to show that such a program is not secure.  We write
$\Sigma \models \alpha\; \ms{secure}^\infty$ for termination-sensitive
noninterference.

\begin{quote}
  We define $\Sigma \models \alpha\; \ms{secure}^\infty$ (program $\alpha$
  satisfies \emph{termination-sensitive noninterference with respect to policy
    $\Sigma$}) \newline iff
  \begin{enumerate}[]
  \item for all $\ell$, $\omega_1$, $\omega_2$, and $\nu_1$,
  \item whenever $\Sigma \vdash \omega_1 \approx_\ell \omega_2$,
  \item and $\omega_1 \lbb \alpha \rbb \nu_1$,
  \item then $\omega_2 \lbb \alpha \rbb \nu_2$ for some $\nu_2$
  \item such that $\Sigma \vdash \nu_1 \approx_\ell \nu_2$.
  \end{enumerate}
\end{quote}

Let's make sure our example program (let's call it $\alpha_0$) is not secure
under this definition.  For this purpose we have to find a counterexample, that
is, two states $\omega_1$ and $\omega_2$ that are low-security equivalent such
that $\alpha_0$ terminates in state $\omega_1$ but not in $\omega_2$.
Fortunately, that is easy: for
\[
  \qquad \alpha_0 = (\mb{if}\; x > 5\; \mb{then}\; (\mb{while}\; \top\; \mb{skip})\; \mb{else}\; \mb{skip})
\]
we can use
\[
  \begin{array}{l@{\hspace{3em}}l}
    \omega_1 = (x \mapsto 0) & \omega_1 \lbb \alpha_0 \rbb (x \mapsto 0) \\
    \omega_2 = (x \mapsto 7) & \mbox{there exists no $\nu_2$ with $\omega_2 \lbb \alpha_0 \rbb \nu_2$}
  \end{array}
\]

\section{Sharpening the Information Flow Type System}

\newcommand{\secureinf}{\ms{secure}^\infty}

The information flow type system so far will admit the program $\alpha_0$ as
secure.  The idea is now to revisit the rules to determine how we might need to
update them to enforce termination-sensitive noninterference.
To parallel the definition of noninterference, we write $\Sigma \vdash
\alpha\; \secureinf$.

The first thing we note is that the security level of expressions and formulas
do not change---they remain the least upper bound of the levels of all variables
occurring in them.  This is due to our decision that expressions and Boolean
conditions always terminate (and are always safe).

So we can focus our attention on the program constructs.  In the fragment we
treat, every command will be safe, that is, we handle \textsc{SafeTiny}.

\paragraph{Assignment.}  We first show the prior (termination-insensitive)
version of the rule.
\begin{rules}
  \infer[\mb{:=}F]
  {\Sigma \vdash x := e\; \ms{secure}}
  {\Sigma \vdash e : \ell
    & \ell' = \Sigma(\mi{pc}) \sqcup \ell
    & \ell' \sqsubseteq \Sigma(x)}
\end{rules}
Because expressions are not involved with nontermination, this rule remains
unchanged!  This is not a proof, of course, it is not difficult to update the
one from the termination-insensitive case.  If $x := e$ terminates in $\omega_1$
then it also will in $\omega_2$, so the two formulations are equivalent in this
case.
\begin{rules}
  \infer[\mb{:=}F^\infty]
  {\Sigma \vdash x := e\; \ms{secure}^\infty}
  {\Sigma \vdash e : \ell
    & \ell' = \Sigma(\mi{pc}) \sqcup \ell
    & \ell' \sqsubseteq \Sigma(x)}
\end{rules}

\paragraph{Sequential Composition.}  Our previous rule can be summarized
as saying that termination-insensitive information flow is \emph{compositional}.
\begin{rules}
  \infer[\mb{\semi}F]
  {\Sigma \vdash \alpha \semi \beta\; \ms{secure}}
  {\Sigma \vdash \alpha\; \ms{secure}
    & \Sigma \vdash \beta\; \ms{secure}}
\end{rules}
Is that still the case, that is, is termination-sensitive information
flow also compositional?  It is not entirely obvious when $\alpha$
does not terminate, but we conjecture:
\begin{rules}
  \infer[\mb{\semi}F^\infty]
  {\Sigma \vdash \alpha \semi \beta\; \ms{secure}^\infty}
  {\Sigma \vdash \alpha\; \ms{secure}^\infty
    & \Sigma \vdash \beta\; \ms{secure}^\infty}
\end{rules}
Let's try to prove or refute the soundness of this.  So we set up:
\begin{tabbing}
  $\Sigma \models \alpha\; \secureinf$ \` (1, first premise) \\
  $\Sigma \models \beta\; \secureinf$ \` (2, second premise) \\
  $\Sigma \vdash \omega_1 \approx_\ell \omega_2$ \` (3, assumption) \\
  $\omega_1 \lbb \alpha \semi \beta \rbb \nu_1$ \` (4, assumption) \\
  $\ldots$ \\
  $\omega_2 \lbb \alpha \semi \beta \rbb \nu_2$ for some $\nu_2$ \` (to show) \\
  with $\Sigma \vdash \nu_1 \approx_\ell \nu_2$ \` (to show)
\end{tabbing}
We start by expanding the definition of the relational semantics for sequential
composition.  We show the new lines in {\color{blue}blue} after each
step.
\begin{tabbing}
  $\Sigma \models \alpha\; \secureinf$ \` (1, first premise) \\
  $\Sigma \models \beta\; \secureinf$ \` (2, second premise) \\
  $\Sigma \vdash \omega_1 \approx_\ell \omega_2$ \` (3, assumption) \\
  $\omega_1 \lbb \alpha \semi \beta \rbb \nu_1$ \` (4, assumption) \\
  \color{blue}$\omega_1 \lbb \alpha \rbb \mu_1$ for some $\mu_1$ \` (5 and) \\
  \color{blue}and $\mu_1 \lbb \beta \rbb \nu_1$ \` (6, from (4) by defn.\ of $\lbb{-}\rbb$) \\
  $\ldots$ \\
  $\omega_2 \lbb \alpha \semi \beta \rbb \nu_2$ for some $\nu_2$ \` (to show) \\
  with $\Sigma \vdash \nu_1 \approx_\ell \nu_2$ \` (to show)
\end{tabbing}
At this point we can use the assumption (coming from the first premise)
that $\alpha$ is secure.
\begin{tabbing}
  $\Sigma \models \alpha\; \secureinf$ \` (1, first premise) \\
  $\Sigma \models \beta\; \secureinf$ \` (2, second premise) \\
  $\Sigma \vdash \omega_1 \approx_\ell \omega_2$ \` (3, assumption) \\
  $\omega_1 \lbb \alpha \semi \beta \rbb \nu_1$ \` (4, assumption) \\
  $\omega_1 \lbb \alpha \rbb \mu_1$ for some $\mu_1$ \` (5, and) \\
  and $\mu_1 \lbb \beta \rbb \nu_1$ \` (6, from (4) by defn.\ of $\lbb{-}\rbb$) \\
  \color{blue}$\omega_2 \lbb \alpha \rbb \mu_2$ for some $\mu_2$ \` (7 and)\\
  \color{blue}with $\Sigma \vdash \mu_1 \approx_\ell \mu_2$ \` (8, from (1), (3), and (5)) \\
  $\ldots$ \\
  $\omega_2 \lbb \alpha \semi \beta \rbb \nu_2$ for some $\nu_2$ \` (to show) \\
  with $\Sigma \vdash \nu_1 \approx_\ell \nu_2$ \` (to show)
\end{tabbing}
At this point we have an evaluation of $\beta$ from a prestate $\mu_1$ and a
state $\mu_2$ that is $\ell$-equivalent to it.  So we can use the security
for $\beta$ (coming from the second premise).
\begin{tabbing}
  $\Sigma \models \alpha\; \secureinf$ \` (1, first premise) \\
  $\Sigma \models \beta\; \secureinf$ \` (2, second premise) \\
  $\Sigma \vdash \omega_1 \approx_\ell \omega_2$ \` (3, assumption) \\
  $\omega_1 \lbb \alpha \semi \beta \rbb \nu_1$ \` (4, assumption) \\
  $\omega_1 \lbb \alpha \rbb \mu_1$ for some $\mu_1$ \` (5, and) \\
  and $\mu_1 \lbb \beta \rbb \nu_1$ \` (6, by defn.\ of $\lbb{-}\rbb$) \\
  $\omega_2 \lbb \alpha \rbb \mu_2$ for some $\mu_2$ (7 and) \\
  with $\Sigma \vdash \mu_1 \approx_\ell \mu_2$ \` (8, from (1), (3), and (5)) \\
  \color{blue}  $\mu_2 \lbb \beta \rbb \nu_2$ for some $\nu_2$ \` (9 and) \\
  \color{blue}  with $\Sigma \vdash \nu_1 \approx_\ell \nu_2$ \` (10, from (2), (6), and (8)) \\
  $\ldots$ \\
  $\omega_2 \lbb \alpha \semi \beta \rbb \nu_2$ for some $\nu_2$ \` (to show) \\
  with $\Sigma \vdash \nu_1 \approx_\ell \nu_2$ \` (to show)
\end{tabbing}
At this point we have evaluations of $\alpha$ and $\beta$ in corresponding
states so we can close the gap simply by composing them.  The $\nu_2$ we
had to find is (not surprisingly) the poststate of $\beta$ when evaluated
in prestate $\mu_2$.
\begin{tabbing}
  $\Sigma \models \alpha\; \secureinf$ \` (1, first premise) \\
  $\Sigma \models \beta\; \secureinf$ \` (2, second premise) \\
  $\Sigma \vdash \omega_1 \approx_\ell \omega_2$ \` (3, assumption) \\
  $\omega_1 \lbb \alpha \semi \beta \rbb \nu_1$ \` (4, assumption) \\
  $\omega_1 \lbb \alpha \rbb \mu_1$ for some $\mu_1$ \` (5, and) \\
  and $\mu_1 \lbb \beta \rbb \nu_1$ \` (6, by defn.\ of $\lbb{-}\rbb$) \\
  $\omega_2 \lbb \alpha \rbb \mu_2$ for some $\mu_2$ (7 and) \\
  with $\Sigma \vdash \mu_1 \approx_\ell \mu_2$ \` (8, from (1), (3), and (5)) \\
  $\mu_2 \lbb \beta \rbb \nu_2$ for some $\nu_2$ \` (9 and) \\
  with $\Sigma \vdash \nu_1 \approx_\ell \nu_2$ \` (10, from (2), (6), and (8)) \\
  $\omega_2 \lbb \alpha \semi \beta \rbb \nu_2$ \` \color{blue}(by defn.\ of $\lbb{-}\rbb$) \\
  with $\Sigma \vdash \nu_1 \approx_\ell \nu_2$ \` \color{blue}(copied from (10))
\end{tabbing}
Good.  Our intuition that even termination-sensitive noninterference
is compositional has been confirmed.

\paragraph{Skip.}  This does nothing, so the rule trivially remains
the same.
\begin{rules}
  \infer[\mb{skip}F^\infty]
  {\Sigma \vdash \mb{skip}\; \ms{secure}^\infty}
  {\mathstrut}
\end{rules}

\paragraph{Conditionals.}  The if-then-else construct somehow doesn't seem to be
involved in nontermination, only loops are.  We therefore expect that the rule
doesn't change.
\begin{rules}
  \infer[\mb{if}F^\infty]
  {\Sigma \vdash \mb{if}\; P\; \mb{then}\; \alpha\; \mb{else}\; \beta\; \ms{secure}^\infty}
  {\Sigma \vdash P : \ell
    & \ell' = \Sigma(\mi{pc}) \sqcup \ell
    & \Sigma' = \Sigma[\mi{pc} \mapsto \ell']
    & \Sigma' \vdash \alpha\; \ms{secure}^\infty
    & \Sigma' \vdash \beta\; \ms{secure}^\infty}
\end{rules}
The proof is just like before, in the termination-insensitive case.  The fact
that we now have to account for the nontermination of $\alpha$ or $\beta$
changes the details of reasoning, of course, but the possibility of
nontermination just comes from the nontermination of each branch.  So
we omit this case.

\paragraph{Loops.}  Here, we'd expect the crux of the changes because loops
introduce nontermination into the language.  First, the earlier rule:
\begin{rules}
  \infer[\mb{while}F]
  {\Sigma \vdash \mb{while}\; P\; \alpha\; \ms{secure}}
  {\Sigma \vdash P : \ell
    & \ell' = \Sigma(\mi{pc}) \sqcup \ell
    & \Sigma' = \Sigma[\mi{pc} \mapsto \ell']
    & \Sigma' \vdash \alpha\; \ms{secure}}
\end{rules}
We see that there are essentially two reasons why a while loop
may leak information using nontermination (with $(x : \ms{H})$)
\begin{enumerate}
\item In the example
  $\mb{if}\; x > 5\; \mb{then}\; (\mb{while}\; \top\; \mb{skip})\; \mb{else}\;
  \mb{skip}$ it is because the loop is in a context where we have
  $\mi{pc} : \ms{H}$.
\item In the example $\mb{while}\; x = 8\; \mb{skip}$ it is because the loop
  guard is of high security.
\end{enumerate}
Since nontermination is observed ``at the outermost level'' and not by a
variable in the local context, it seems we have to require that both the
security level of the $\mi{pc}$ and the guard are $\bot$, that is, the least
element of the lattice.  This means that $\ell' = \Sigma(\mi{pc}) \sqcup \ell$
should also be $\bot$, and we don't need to update the security level of
$\mi{pc}$.
\begin{rules}
  \infer[\mb{while}F^\infty]
  {\Sigma \vdash \mb{while}\; P\; \alpha\; \ms{secure}^\infty}
  {\Sigma \vdash P : \bot
    & \Sigma(\mi{pc}) = \bot
    & \Sigma \vdash \alpha\; \ms{secure}^\infty}
\end{rules}
Is this sufficient to guarantee termination-sensitive noninterference?
We set up:
\begin{tabbing}
  $\Sigma \models \alpha\; \secureinf$ \` (1, premise) \\
  $\Sigma \vdash \omega_1 \approx_\ell \omega_2$ \` (2, assumption) \\
  $\omega_1 \lbb \mb{while}\; P\; \alpha \rbb \nu_1$ \` (3, assumption) \\
  $\ldots$ \\
  $\omega_2 \lbb \mb{while}\; P\; \alpha \rbb \nu_2$ for some $\nu_2$ \` (to show) \\
  with $\Sigma \vdash \nu_1 \approx_\ell \nu_2$ \` (to show)
\end{tabbing}
At this point, there is a small hiccup: the evaluation of a loop may refer back
again to the evaluation of the same loop.  We therefore use an auxiliary
induction on the number of times we iterate the loop when computing the
poststate $\nu_1$ from $\omega_1$.  We know it does
compute a poststate, so the number of steps is finite and the induction makes
sense.  Let's call that number $n$ (which is also the notation we chose in the
semantic definition of $\omega \lbb \mb{while}\; P\; \alpha\rbb^n \nu$, see
\href{\courseurl/lectures/02-semantics.pdf}{Lecture 2}).
\begin{description}
\item[Case:] $n = 0$.  Then $\omega_1 \not\models P$ and
  $\nu_1 = \omega_1$.  Because $P$ only contains variables of security level
  $\bot \sqsubseteq \ell$ and $\Sigma \vdash \omega_1 \approx_\ell \omega_2$, we
  also have $\omega_2 \not\models P$.  (See Lemma 2 in
  \href{\courseurl/lectures/12-infoflow.pdf}{Lecture 12}.)  Therefore
  $\nu_2 = \omega_2$ will satisfy the requirements of the theorem.
\item[Case:] $n > 0$.  Then $\omega_1 \models P$ and
  $\omega_1 \lbb \alpha \semi \mb{while}\; P\; \alpha \rbb \nu_1$.  As in the
  previous case $\omega_2 \models P$ and it remains to show that
  $\omega_2 \lbb \alpha \semi \mb{while}\; P\; \alpha \rbb \nu_2$ for some
  $\nu_2$ with $\Sigma \vdash \nu_1 \approx_\ell \nu_2$.

  As in the case for composition (and exploiting the security of $\alpha$ that
  comes from the premise) this reduces to showing that
  $\mu_1 \lbb \mb{while}\; P\; \alpha \rbb \nu_1$ implies
  $\mu_2 \lbb \mb{while}\; P\; \alpha \rbb \nu_2$ for some $\nu_2$
  indistinguishable from $\nu_1$ at level $\ell$, with the knowledge that
  $\Sigma \vdash \mu_1 \approx_\ell \mu_2$.  This is true because of the
  induction hypothesis on $n-1 < n$, which is the remaining number of times the
  loop is traversed.
\end{description}

\paragraph{Tests.}  Recall that $\mb{test}\; P$ aborts if $P$ is false.  This
means it represents another type of program (besides a loop) that does not have a
poststate.  Lumping together aborting a program and nontermination may be
questionable in practice, but if we consider ``nontermination'' simply as
representing the absence of a poststate, we would have defined
\[
  \mb{test}\; P \triangleq \mb{while}\; \lnot P\; \mb{skip}
\]
Then the rule derived from these considerations would be:
\begin{rules}
  \infer[\mb{test}F^\infty]
  {\Sigma \vdash \mb{test}\; P\; \ms{secure}^\infty}
  {\Sigma \vdash P : \bot & \Sigma(\mi{pc}) = \bot}
\end{rules}

All together we obtain soundness.

\begin{theorem}[Soundness of termination-sensitive information flow types]
  If $\Sigma \vdash \alpha\; \secureinf$ then $\Sigma \models \alpha\; \secureinf$
\end{theorem}
\begin{proof}
  All the rules are \emph{sound}, that is, the preserve semantic validity of the
  premises.  By a trivial induction over the derivation of $\Sigma \vdash \alpha\; \secureinf$,
  every program judged secure satisfies termination-sensitive noninterference.
\end{proof}

\section{Further Discussion}

The soundness of all rules for $\Sigma \vdash \alpha\; \secureinf$ guarantees
that it implies $\Sigma \models \alpha\; \secureinf$, that is,
termination-sensitive information flow.  The other direction (completeness) is
not true, and there are simple counterexamples.

We also have the intuition that it should be \emph{stricter} than the
termination-insensitive one.  That should be true syntactically, in the type system,
and semantically, with respect to noninterference.  Let's check that:

\begin{theorem}
  If $\Sigma \models \alpha\; \secureinf$ then $\Sigma \models \alpha\; \ms{secure}$.
\end{theorem}
\begin{proof}
  We set up:
  \begin{tabbing}
    $\Sigma \models \alpha\; \secureinf$ \` (1, assumption) \\
    $\Sigma \vdash \omega_1 \approx_\ell \omega_2$ \` (2, assumption) \\
    $\omega_1 \lbb \alpha \rbb \nu_1$ \` (3, assumption) \\
    $\omega_2 \lbb \alpha \rbb \nu_2$ \` (4, assumption) \\
    $\ldots$ \\
    $\Sigma \vdash \nu_1 \approx_\ell \nu_2$ \` (to show)
  \end{tabbing}
  We can now apply the definition of $\secureinf$.
  \begin{tabbing}
    $\Sigma \models \alpha\; \secureinf$ \` (1, assumption) \\
    $\Sigma \vdash \omega_1 \approx_\ell \omega_2$ \` (2, assumption) \\
    $\omega_1 \lbb \alpha \rbb \nu_1$ \` (3, assumption) \\
    $\omega_2 \lbb \alpha \rbb \nu_2$ \` (4, assumption) \\
    \color{blue}$\omega_2 \lbb \alpha \rbb \nu_2'$ for some $\nu_2'$ with \` (5 and) \\
    \color{blue}$\Sigma \vdash \nu_1 \approx_\ell \nu_2'$ \` (6, from (1), (2), and (3)) \\
    $\ldots$ \\
    $\Sigma \vdash \nu_1 \approx_\ell \nu_2$ \` (to show)
  \end{tabbing}
  Since our language is deterministic, we have $\nu_2 = \nu_2'$ and the
  proof is complete.
  \begin{tabbing}
    $\Sigma \models \alpha\; \secureinf$ \` (1, assumption) \\
    $\Sigma \vdash \omega_1 \approx_\ell \omega_2$ \` (2, assumption) \\
    $\omega_1 \lbb \alpha \rbb \nu_1$ \` (3, assumption) \\
    $\omega_2 \lbb \alpha \rbb \nu_2$ \` (4, assumption) \\
    \color{blue}$\omega_2 \lbb \alpha \rbb \nu_2'$ for some $\nu_2'$ with \` (5 and) \\
    \color{blue}$\Sigma \vdash \nu_1 \approx_\ell \nu_2'$ \` (6, from (1), (2), and (3)) \\
    $\Sigma \vdash \nu_1 \approx_\ell \nu_2$ \` \color{blue}(7, from (4), (5), and (6) by determinism)
  \end{tabbing}
\end{proof}

Any syntactic (decidable) type system for a Turing-complete language may be
considered an approximation of a true semantic property.  This is the essence of
Rice's theorem.  Any type system that restricts information flow then represents
a tradeoff between the kind of programs that are allowed, and the simplicity and
uniformity of the system.  This is difficult to discuss in the abstract, so in
Lab 2 you will have the opportunity to explore it a bit yourself.  The
information flow type systems from the last three lectures all seems to
represent reasonable compromises: the rules are sound and quite simple, and many
programs can be written (even if some require declassification).

Perhaps the most significant issue is that there exist further \emph{side
  channels} by which high-security information can be obtained, most notably by
observing program runtime.  We will complete this investigation in the next
lecture.

\section{Practice Problems}

These problems are for self-study and are not to be handed in.

\subsection{Classifying Programs [3 parts]}

Consider the two-level security lattice with $\ms{L} \sqsubset \ms{H}$ and
the following security policy:
\[
  \Sigma = (x : \ms{H},\; y : \ms{L},\; z : \ms{L})
\]
For each of the following programs, determine whether it satisfies (a)
termination-insensitive noninterference ($\Sigma \models \alpha\; \ms{secure}$),
(b) termination-sensitive noninterference
($\Sigma \models \alpha\; \secureinf$), and (c) whether it can be derived in the
corresponding type system ($\Sigma \vdash \alpha\; \ms{secure}$ or
$\Sigma \vdash \alpha\; \secureinf$).  Justify each answer with either a
derivation, a brief argument, or a counterexample as appropriate.

\begin{task}\rm
  \label{task:p1a}
  $\alpha_1 = (\mb{while}\; y > 0\; (y := y - 1))$
\end{task}

\begin{task}\rm
  \label{task:p1b}
  $\alpha_2 = (\mb{while}\; y > 0\; (y := y - x))$
\end{task}

\begin{task}\rm
  \label{task:p1c}
  $\alpha_3 = (\mb{if}\; x > 0\; \mb{then}\; (\mb{while}\; y > 0\; (y := y - 1))\; \mb{else}\; z := 0)$
\end{task}

\subsection{Bounded Loops [4 parts]}

We add a new construct $\mb{repeat}\; e\; \alpha$ to our language, where $e$ is
an integer expression.  Its intended behavior is to execute the body $\alpha$
exactly $\omega \lbb e \rbb$ times when $\omega \lbb e \rbb \geq 0$, and to
behave as $\mb{skip}$ when $\omega \lbb e \rbb < 0$.

For example, if $\omega = (n \mapsto 3, y \mapsto 0)$ then
\[
  \omega \lbb \mb{repeat}\; n\; (y := y + 1) \rbb (n \mapsto 3, y \mapsto 3)
\]
Note that unlike $\mb{while}$, the bound $e$ is evaluated \emph{once} in the
initial state, and $\mb{repeat}$ always terminates.

\begin{task}\rm
  \label{task:p2a}
  Give a semantic definition of
  $\omega \lbb \mb{repeat}\; e\; \alpha \rbb \nu$.  You may use an auxiliary
  definition indexed by the number of remaining iterations.
\end{task}

\begin{task}\rm
  \label{task:p2b}
  Propose a rule for $\mb{repeat}$ in the termination-sensitive information flow
  type system:
  \[
    \infer[\mb{repeat}F^\infty]
    {\Sigma \vdash \mb{repeat}\; e\; \alpha\; \secureinf}
    {\ldots}
  \]
  Your rule should be as permissive as possible while remaining sound.
  In particular, consider whether $\mb{repeat}$ requires the same restrictions
  on $e$ and $\Sigma(\mi{pc})$ as $\mb{while}$.
\end{task}

\begin{task}\rm
  \label{task:p2c}
  Give a brief argument for the soundness of your rule, following the structure
  of the while loop case from this lecture.  You do not need to write a
  complete formal proof, but you should identify the key property of
  $\mb{repeat}$ that makes the argument go through.
\end{task}

\begin{task}\rm
  \label{task:p2d}
  Give a concrete program using $\mb{repeat}$ and a security policy $\Sigma$
  such that $\Sigma \vdash \alpha\; \secureinf$ is derivable using your rule
  from \autoref{task:p2b}, but would \textbf{not} be derivable if
  $\mb{repeat}\; e\; \alpha$ were desugared to
  $i := 0 \semi \mb{while}\; (i < e)\; (\alpha \semi i := i + 1)$ using the
  $\mb{while}F^\infty$ rule.  Briefly explain why.
\end{task}

\bibliographystyle{plainnat}
\bibliography{bibliography}

\end{document}
